\documentclass[11pt]{article}
\usepackage[a4paper,margin=25mm]{geometry}
\usepackage[T1]{fontenc}
\usepackage{lmodern,amsmath,amssymb,amsthm,booktabs,array}
\usepackage[hidelinks]{hyperref}
\newtheorem{theorem}{Theorem}
\newtheorem{corollary}[theorem]{Corollary}
\newcommand{\di}{d_I}
\newcommand{\Z}{\mathbb Z}
\newcommand{\Q}{\mathbb Q}
\newcommand{\coef}[1]{[z^{#1}]}
\title{A general-deficit extension of Shao and Zhao's\\Italian cycle-domination formulas}
\author{Haobo Ma\thanks{ChronoAI Pte Ltd, Singapore.}\and
Wenlin Zhang\thanks{National University of Singapore; ChronoAI Pte Ltd, Singapore.}}
\date{17 November 2026}
\begin{document}
\maketitle
\begin{abstract}
By a gap decomposition, we extend Shao and Zhao's Theorem~7.1 (2026, arXiv:\allowbreak2610.00108v1,
pp.~22--24) to every deficit, answering the fourth future-work direction
in their Section~9, p.~29: a finite binomial sum counts the Italian
dominating functions of weight $2n-j$ on a labelled cycle for every
$n\geq3$ and $0\leq j\leq2n$.
The proof cuts the deficit word at its zeros and separates isolated twos
from the exceptional gap $22$; marking a vertex handles periodic words
without passing to rotation classes.
The same decomposition gives a rational bivariate generating function,
the total-count recurrence asked for in the second direction, and an
explicit stable range for the polynomiality of each fixed-deficit coefficient.
We also determine the degree and leading coefficient of the exclusion
polynomial and prove that paths have the same total-count growth ratio.
\end{abstract}

\section{The coefficient formula}
Let $C_n$ have labelled vertices $\Z/n\Z$, where $n\geq3$.
An Italian dominating function is a map $f:\Z/n\Z\to\{0,1,2\}$ such that
$f(i)=0$ implies $f(i-1)+f(i+1)\geq2$. Its weight is $\sum_i f(i)$;
$\di(C_n,k)$ counts the functions of weight $k$, and is zero outside
$0\leq k\leq2n$. No identification under rotation or reflection is made.
Put $y_i=2-f(i)$ and call $j=\sum_i y_i$ the deficit. Then
\begin{equation}\label{eq:local}
 y_i=2\ \Longrightarrow\ y_{i-1}+y_{i+1}\leq2.
\end{equation}
This is Shao and Zhao's Corollary~3.2, p.~7~\cite{SZ}.
Their boundary-conditioned dynamic program computes all cycle coefficients;
Theorem~7.1 gives explicit formulas for $j\leq5$, while Section~9 asks for
general $j$. Here the answer is a positive finite sum rather than a
case-by-case subtraction of forbidden configurations. Local restrictions
on compositions also have a general transfer-matrix treatment~\cite{BC};
the particular gap classification below makes every coefficient explicit.

For integers $u,v$, set $\binom uv=0$ unless $0\leq v\leq u$, and use the
ordinary binomial coefficient otherwise. Define $r=j-n$, $s=2n-j$, and
\[
 Z_{n,j}=\begin{cases}\displaystyle\frac ns\binom sr,
 &r\geq0,\ 2r\leq n,\ s>0,\\0,&\text{otherwise},\end{cases}
 \qquad
 E_{n,j}=\begin{cases}3,&3\mid n\text{ and }3j=4n,\\0,&\text{otherwise}.
 \end{cases}
\]
For nonnegative integers $a,b$, write
\[
 c=j-2a-4b,\qquad t=n-j+a+b,\qquad m=t+b,
\]
and define
\begin{equation}\label{eq:term}
 G_{n,j}(a,b)=
 \begin{cases}
 \displaystyle\frac nm\binom mb\binom{c+t-1}{c}\binom{c+t}{a},
 &c\geq0\text{ and }t\geq1,\\[3pt]
 0,&\text{otherwise}.
 \end{cases}
\end{equation}
The positive branch has $m\geq1$, so its denominator is nonzero.

\begin{theorem}\label{thm:main}
For $n\geq3$ and $0\leq j\leq2n$,
\begin{equation}\label{eq:main}
 \di(C_n,2n-j)=Z_{n,j}+E_{n,j}
       +\sum_{b=0}^{j}\sum_{a=0}^{j}G_{n,j}(a,b).
\end{equation}
\end{theorem}

\begin{proof}
We count words $y$ satisfying~\eqref{eq:local} and having sum $j$.
First suppose $y$ contains no zero. Its letters are $1$ and $2$;
two adjacent twos would violate~\eqref{eq:local}. There are $r=j-n$
twos and $s=2n-j$ ones. Thus $r\geq0$, $r\leq s$, and $s>0$.
Mark one of the ones and read clockwise from it. Among the $s$ slots
immediately following the ones, choose $r$ slots in which to place a two.
There are $\binom sr$ rooted arrangements and $n$ choices for the
label of the marked one. Each unmarked labelled word has exactly $s$
marks, giving $Z_{n,j}$. This argument counts actual marked vertices,
including those in periodic words.

Now suppose there are $m\geq1$ zeros. The $m$ successive gaps between
zeros are linear words on $\{1,2\}$, with the empty word allowed.
If a gap contains adjacent twos, the two outer neighbors of this pair
must both be zero by~\eqref{eq:local}. Consequently that entire gap
is exactly $22$. Call it exceptional. Every other gap has no adjacent
twos and will be called ordinary. Conversely, all these gaps satisfy
the local rule, including at the zero boundaries.

Let $b$ be the number of exceptional gaps, $t=m-b$ the number of ordinary
gaps, $c$ the total number of ones, and $a$ the number of isolated twos
in ordinary gaps. Counting length and deficit gives
\begin{equation}\label{eq:inventory}
 n=m+c+a+2b,\qquad j=c+2a+4b.
\end{equation}
If $t=0$, the whole word repeats $022$: necessarily $3\mid n$ and
$3j=4n$, and the three choices for the residue class of the zeros are
the only possibilities. This is $E_{n,j}$.

Suppose $t\geq1$. Mark a zero and list the $m$ gaps clockwise from it.
Choose the $b$ exceptional gaps in $\binom mb$ ways. Distribute the
$c$ ones among the $t$ ordinary gaps by a weak composition, in
$\binom{c+t-1}{c}$ ways. If a particular ordinary gap contains $c_i$
ones, its isolated twos may occupy any subset of its $c_i+1$ slots
(before, between, or after those ones). Across the ordinary gaps there
are $c+t$ slots. Choosing $a$ of these gives $\binom{c+t}{a}$ choices.
The construction is bijective, because deleting the isolated twos
recovers the distributed ones and the chosen slots.

There are $n$ choices for the label of the marked zero, while each
unmarked word has exactly $m$ possible marks. Hence its contribution is
the positive branch of~\eqref{eq:term}. Solving~\eqref{eq:inventory}
gives precisely $c,t,m$ defined above. Every contributing word has
$0\leq a,b\leq j$; impossible inventories have zero contribution.
The three classes are disjoint and exhaustive, proving~\eqref{eq:main}.
In particular $j=0$ gives the unique all-zero deficit word, and $j=2n$
gives zero, since the all-two deficit word violates the local rule.
\end{proof}

\section{Generating function and consequences}
Let
\[
 A_n(z)=\sum_{j=0}^{2n}\di(C_n,2n-j)z^j \quad(n\geq3),
\]
and let $x$ mark length. All generating functions below are formal power
series in $x$ with coefficients in $\Q[z]$. Set
\begin{align*}
 D(x,z)&=1-xz-x^2z^3,\\
 H(x,z)&=x\left(\frac{1+xz^2}{D(x,z)}+x^2z^4\right),\\
 P(x,z)&=D(x,z)(1-H(x,z))\\
 &=1-(1+z)x-(z^2+z^3)x^2-z^4x^3+z^5x^4+z^7x^5.
\end{align*}
The factor $x$ in $H$ counts a zero and the following gap.

\begin{corollary}\label{cor:gen}
For every $n\geq3$,
\begin{equation}\label{eq:log}
 A_n(z)=[x^n]\left(-x\frac{\partial}{\partial x}\log P(x,z)\right).
\end{equation}
Consequently, for $n\geq3$,
\begin{align}\label{eq:weighted}
 A_{n+5}(z)={}&(1+z)A_{n+4}(z)+(z^2+z^3)A_{n+3}(z)\notag\\
 &+z^4A_{n+2}(z)-z^5A_{n+1}(z)-z^7A_n(z).
\end{align}
Writing $a_n=A_n(1)$ gives
\[
 a_{n+5}=2a_{n+4}+2a_{n+3}+a_{n+2}-a_{n+1}-a_n,
 \qquad (a_3,a_4,a_5,a_6,a_7)=(23,60,167,467,1297).
\]
\end{corollary}
\begin{proof}
An ordinary gap with $c$ ones has $c+1$ slots for isolated twos, so its
generating function is
\[
 \sum_{c\geq0}(xz)^c(1+xz^2)^{c+1}
       =\frac{1+xz^2}{1-xz-x^2z^3}.
\]
The exceptional gap contributes $x^2z^4$. Thus the same zero-marking
argument as in Theorem~\ref{thm:main} gives the zero-containing contribution
\[
 x\partial_x\sum_{m\geq1}\frac{H(x,z)^m}{m}
       =-x\partial_x\log(1-H(x,z)).
\]
Indeed differentiation supplies the factor $n$, and division by $m$
removes the zero mark without using any orbit-size assumption.
For a zero-free word mark a one instead. A one and its optional following
two have generating function $B=xz(1+xz^2)$, so the identical argument
gives $-x\partial_x\log(1-B)=-x\partial_x\log D$.
Every zero-free admissible cycle of length at least three contains a one.
Adding the two contributions proves~\eqref{eq:log}. The formal expression
also has coefficients at lengths one and two, but we give those no
simple-cycle interpretation.

Multiplying $-xP_x/P$ by $P$ shows its coefficients satisfy the displayed
order-five recurrence at degrees greater than five. Taking degree $n+5$
with $n\geq3$ uses only simple-cycle coefficients. Setting $z=1$ gives
the total-count recurrence. The initial totals follow by substituting
$n=3,4,5,6,7$ in~\eqref{eq:main} and summing over $j$.
\end{proof}

\begin{corollary}\label{cor:poly}
For each fixed $j\geq0$, $\di(C_n,2n-j)$ agrees, for
$n\geq\max(3,j+1)$, with a polynomial in $n$ of degree $j$ and leading
coefficient $1/j!$.
\end{corollary}
\begin{proof}
When $n>j$, neither $Z_{n,j}$ nor $E_{n,j}$ contributes, and $t\geq1$
for every summand with $c\geq0$. The finite set of pairs $a,b$ is
independent of $n$. If $b\geq1$, the factor
$\binom mb/m=\binom{m-1}{b-1}/b$ is polynomial in $n$.
If $b=0$ and $c\geq1$, $\binom{c+m-1}{c}$ contains a factor $m$.
If $b=c=0$ and $a\geq1$, $\binom ma$ contains a factor $m$ instead.
Thus every summand is polynomial of degree at most
$c+a+b=j-a-3b$. The pair $a=b=0$ contributes $\binom nj$:
for $j>0$ this follows from
$n\binom{n-1}{j}/(n-j)=\binom nj$, and for $j=0$ it is $1$.
All other pairs have smaller degree. This proves the assertion and
supplies a concrete sufficient stable range for Lemma~3.4 of~\cite{SZ}.
\end{proof}

\begin{corollary}\label{cor:excluded}
Let $T(n,j)=[z^j](1+z+z^2)^n$ and
$F_j(n)=T(n,j)-\di(C_n,2n-j)$.
For $j\leq4$, $F_j(n)=0$ whenever $n\geq3$ and $j\leq2n$.
For every $j\geq5$ and $n\geq j+1$, $F_j(n)$ is a polynomial in $n$
of degree exactly $j-4$, with leading coefficient $2/(j-5)!$.
\end{corollary}
\begin{proof}
A forbidden triple is $122$, $221$, or $222$, so none has deficit below
five. For fixed $j$ and $n>j$, every word of deficit $j$ has a zero.
Decompose an excluded word into its $h\geq1$ nonempty positive runs,
separated by nonempty zero gaps. At least one run contains a forbidden
triple and has weight at least five; all other runs have weight at least
one. Consequently $h\leq j-4$.

Fix an ordered list of these runs, of total length $\ell$ and weight $j$.
There are finitely many such lists, independently of $n$.
Mark the start of one run. Assign its vertex label and distribute
the $n-\ell$ zeros among the $h$ nonempty gaps. Summing
\[
 \frac nh\binom{n-\ell-1}{h-1}
\]
over the ordered run lists counts the unmarked labelled words, since
each word has exactly $h$ marks. Each summand is a polynomial of
degree $h$ for $n>j$, including its zero values when there are too few
zeros for the gaps. Thus $\deg F_j\leq j-4$.
Equality $h=j-4$ requires one weight-five bad run, necessarily $122$ or
$221$, and $j-5$ singleton runs $1$. There are $2h$ such ordered lists,
all with $\ell=j-2$. Their combined contribution is
$2n\binom{n-j+1}{j-5}$; all other lists give smaller degree.
This proves the stated degree and leading coefficient, strengthening
the degree bound in Lemma~3.4 of~\cite{SZ} with a sufficient stable range.
\end{proof}

The low deficits recover Theorem~7.1 of~\cite{SZ} without an expansion:
put $T(n,j)=\coef{j}(1+z+z^2)^n$. The smallest forbidden triples in
\eqref{eq:local} are $122,221,222$, with deficits $5,5,6$.
Therefore $\di(C_n,2n-j)=T(n,j)$ for $0\leq j\leq4$.
For $j=5$ and $n\geq4$, the invalid words have two adjacent twos and
one adjacent one; the $n$ choices of the pair and its two ends give
$\di(C_n,2n-5)=T(n,5)-2n$. For $n=3$ that coefficient is zero instead.
These are recoveries of the cited formulas, not additional cases of a
new low-deficit theorem.
For $j\leq3$, the corresponding Roman formulas were proved earlier by
Deepak, Alwardi and Indiramma (2021), Theorem~2.4(ix)--(xii), p.~91~\cite{DAI};
Corollary~\ref{cor:roman} explains why they also count Italian functions.

\begin{corollary}\label{cor:growth}
The total cycle count satisfies $a_n\sim\lambda^n$ and
$a_{n+1}/a_n\to\lambda$, where $2.784<\lambda<2.785$ is the unique
root of largest modulus of
\[
 q(u)=u^5-2u^4-2u^3-u^2+u+1.
\]
\end{corollary}
\begin{proof}
At $z=1$, the positive-coefficient series
$H(x,1)=x(1+x)/(1-x-x^2)+x^3$ increases from zero to infinity as
$x$ increases from zero to $\alpha=(\sqrt5-1)/2$.
There is a unique $\rho\in(0,\alpha)$ with $H(\rho,1)=1$.
For $|x|<\rho$, $|H(x,1)|\leq H(|x|,1)<1$, and $D(x,1)\ne0$.
Thus $P(x,1)$ has no zero in that disk. On $|x|=\rho$, equality
in the triangle inequality and the nonzero coefficients of $x$ and
$x^2$ force any zero of $1-H(x,1)$ to be $x=\rho$.
This zero is simple, since $H_x(\rho,1)>0$ and $D(\rho,1)>0$.
Writing $P(x,1)=\prod_{i=1}^5(1-\lambda_i x)$, with roots repeated
according to multiplicity, equation~\eqref{eq:log} gives
$a_n=\sum_i\lambda_i^n$ for $n\geq3$.
Exactly one $\lambda_i=1/\rho$ has largest modulus, proving both limits.
It is the unique positive root greater than $1/\alpha$; the exact signs
$q(2.784)<0<q(2.785)$ locate it as stated. This establishes the
cycle-growth observation in Section~8.5(8) of~\cite{SZ}.
\end{proof}

\begin{corollary}\label{cor:path}
For $p_n=\sum_k\di(P_n,k)$ there is a constant $C>0$ such that
$p_n\sim C\lambda^n$. In particular $p_{n+1}/p_n\to\lambda$, with the
same $\lambda$ as in Corollary~\ref{cor:growth}.
\end{corollary}
\begin{proof}
Index a matrix $M$ by the nine pairs $(a,b)\in\{0,1,2\}^2$.
Its entry from $(a,b)$ to $(b,c)$ is one when $b\ne0$ or $a+c\geq2$;
all other entries are zero. Closed walks give
$a_n=\operatorname{tr}(M^n)$ for $n\geq3$: the consecutive pairs of
a closed walk recover exactly the original labelled function.
For any two states there is a four-step walk through $(2,2)$, namely
\[
 (a,b),\ (b,2),\ (2,2),\ (2,c),\ (c,d).
\]
Every displayed transition is allowed, so $M^4$ is strictly positive.
The finite Perron--Frobenius theorem therefore gives a simple dominant
eigenvalue $\mu>0$ with positive left and right eigenvectors, and
$\operatorname{tr}(M^n)\sim\mu^n$. Corollary~\ref{cor:growth} identifies
$\mu=\lambda$.
For $n\geq2$, $p_n=u^{\mathsf T}M^{n-2}v$, where
$u_{(a,b)}=1$ iff $a\ne0$ or $b=2$, and
$v_{(a,b)}=1$ iff $b\ne0$ or $a=2$.
These are precisely the two endpoint constraints. Both vectors are
nonnegative and nonzero, so their pairings with the positive Perron
eigenvectors are positive. The same theorem gives the claimed
asymptotic and ratio. This proves the path part of Section~8.5(8)--(9)
of~\cite{SZ}, not just the cycle part.
\end{proof}

The next identities also follow in their common ranges by subtracting
the formulas in Theorem~2.4(ix)--(xiv) of~\cite{DAI} from Theorem~7.1
of~\cite{SZ}. The local proof includes the small-cycle boundaries.
\begin{corollary}\label{cor:roman}
Let $r(C_n,k)$ count Roman dominating functions: each zero must have
a neighbor labelled two. For $n\geq3$ the difference
$\di(C_n,2n-j)-r(C_n,2n-j)$ is zero for $0\leq j\leq3$, is $n$ for
$j=4$, and is $n(n-3)$ for $j=5$.
\end{corollary}
\begin{proof}
An Italian function fails the Roman condition exactly when some zero
has both neighbors one. In deficit variables this is $121$, of weight
four. At deficit four its center has $n$ choices and the remaining
vertices are zero. At deficit five there is one further deficit-one
vertex, in any of the $n-3$ positions outside that triple. In either case
the center is the unique deficit-two vertex, so there is no multiple
counting. These words satisfy the Italian condition. The argument
also covers $n=3$, when no further position is available. This makes
the comparison of Remark~7.4 of~\cite{SZ} quantitative.
\end{proof}

\begin{corollary}\label{cor:star}
For a star with $m\geq2$ leaves, its weight polynomial is
\[
 D_I(K_{1,m},u)=u^2(1+u+u^2)^m+(1+u)(u+u^2)^m.
\]
Consequently $\di(K_{1,m},2)=1$ for $m\geq3$, whereas the values
for $m=1,2$ are respectively $3,2$.
\end{corollary}
\begin{proof}
A center labelled two permits arbitrary leaf labels, giving the first
term. A center labelled one forces all leaves positive, giving
$u(u+u^2)^m$. A zero center also forces positive leaves; for $m\geq2$
their sum is automatically at least two, giving $(u+u^2)^m$.
Reading the coefficient of $u^2$ proves the assertion for $m\geq2$.
For $m=1$ the weight-two functions are $(0,2),(1,1),(2,0)$.
This proves the observation in Section~8.5(10) of~\cite{SZ}.
\end{proof}

\section{Coverage and verification scope}
The first six rows address Shao and Zhao's future-work bullets in
Section~9, p.~29~\cite{SZ}. An answered request is distinguished from
a previously stated theorem or a numerical observation.
\begin{center}\small
\begin{tabular}{@{}p{.46\linewidth}p{.46\linewidth}@{}}
\toprule
Source item & Coverage \\
\midrule
Broader graph classes (bullet 1) & Not covered.\\
Total cycle-count recurrence (bullet 2) & Answered: yes;
Corollary~\ref{cor:gen}.\\
Unimodality (bullet 3; also \S8.5(5)) & Not covered.\\
General-deficit cycle formulas (bullet 4) & Answered: yes;
Theorem~\ref{thm:main}, by a positive gap sum.\\
$q$-analogues and Roman relations (bullet 5) & Roman low-deficit
comparison: Corollary~\ref{cor:roman}; $q$-analogues not covered.\\
Graphs sharing the same polynomial (bullet 6) & Not covered.\\
\midrule
Theorem~7.1; \S8.5(6), low deficits & Reproved; prior proof:
Shao and Zhao (2026); the paragraph after Corollary~\ref{cor:excluded}.\\
Lemma~3.4, eventual polynomiality and exclusion degree & Reproved;
prior proof: Shao and Zhao (2026). Explicit range and exact exclusion
degree/leading term: Corollaries~\ref{cor:poly}--\ref{cor:excluded}.\\
Remark~7.4, Italian--Roman comparison & Reproved from the formulas of
Deepak et al. (2021) and Shao--Zhao (2026): Corollary~\ref{cor:roman}.\\
\S8.5(7)--(9), path/cycle totals and growth & Proved:
Corollaries~\ref{cor:growth}--\ref{cor:path}; finite counts in \texttt{checks.json}.\\
\S8.5(10), star coefficient & Proved: Corollary~\ref{cor:star}.\\
\S8.5(1)--(4), endpoint weights and dominance & Prior results cited:
Propositions~2.1--2.3 and Remark~7.1 of Shao and Zhao (2026).\\
\S8.4, other specific trees & Not covered.\\
\bottomrule
\end{tabular}
\end{center}

The Lean declaration \texttt{CycleDeficit.result} proves exactly~\eqref{eq:main} for $n\geq3$
and $j\leq2n$, with its functions \texttt{zeroFree},
\texttt{allExceptional}, and \texttt{gapTerm} corresponding to $Z,E,G$.
The underlying \texttt{dI} counts labelled functions of their given weight.
Its source is the pinned
\href{https://github.com/the-omega-institute/trureturing/blob/7d127fe3daacaa59991ac9a3ed5e2931a1eeadeb/D5/S3/Combinatorics/Graph/ItalianDomination/CycleDeficit.lean}{\texttt{CycleDeficit.lean}}.
The total recurrence and its five initial values are covered separately by
\href{https://github.com/the-omega-institute/trureturing/blob/7d127fe3daacaa59991ac9a3ed5e2931a1eeadeb/D5/S3/Combinatorics/Graph/CycleItalianDominationRecurrence.lean}{\texttt{result} in the cycle-recurrence source}.
The auxiliary source
\href{https://github.com/the-omega-institute/trureturing/blob/7d127fe3daacaa59991ac9a3ed5e2931a1eeadeb/D5/S3/Combinatorics/Graph/ItalianDomination/OrdinaryWords.lean}{\texttt{OrdinaryWords.lean}}
contains \texttt{dI\_eq\_deficit\_card} (the weight--deficit correspondence)
and \texttt{ordinaryFamily\_card} (ordinary gap families).
The source
\href{https://github.com/the-omega-institute/trureturing/blob/7d127fe3daacaa59991ac9a3ed5e2931a1eeadeb/D5/S3/Combinatorics/Graph/ItalianDomination/RootedCycleGaps.lean}{\texttt{RootedCycleGaps.lean}}
contains \texttt{refined\_cycle\_marking\_count}, the marked-zero counting
identity.

The generating function, weighted recurrence, stable range, exclusion leading
term, path/cycle growth, Roman differences and star formula have the written
proofs above; these are distinct from the two main Lean statements.

The accompanying \texttt{verify.py} compares the coefficient formula with
an independent count with fixed first two labels, and with direct enumeration of
functions on small cycles. It also checks finite instances of the recurrence
and polynomial differences. Exact ranges and outputs are in
\texttt{checks.json}; these finite checks complement the all-parameter proofs.

\begin{thebibliography}{9}
\bibitem{SZ} P. Shao and C. Zhao,
\emph{Counting Weight-$k$ Italian Dominating Sets on Trees and Cycles},
arXiv:2610.00108v1 (2026).
\url{https://arxiv.org/abs/2610.00108v1}.
\bibitem{DAI} G. Deepak, A. Alwardi and M. H. Indiramma,
Roman domination polynomial of cycles,
\emph{International Journal of Mathematical Combinatorics}
\textbf{3} (2021), 86--97.
\url{https://fs.unm.edu/IJMC/IJMC-3-2021.pdf}.
\bibitem{BC} E. A. Bender and E. R. Canfield,
Locally restricted compositions II: General restrictions and infinite matrices,
\emph{Electronic Journal of Combinatorics} \textbf{16}(1) (2009), R108.
\url{https://doi.org/10.37236/197}.
\end{thebibliography}
\bigskip
\noindent\small This note is licensed under Creative Commons Attribution 4.0
International (CC BY 4.0).
\end{document}
